\chapter{I provider di deployment}
\label{cap:15-deployment-provider}

Due moduli implementano l'interfaccia \code{ManagedDeploymentProvider}
(\cref{cap:07-modalita-esecuzione}) e forniscono alla modalit\`a
\code{DEPLOYMENT} un meccanismo concreto: \code{k8s-deployment-provider}
(1\,202 righe) e \code{container-deployment-provider} (1\,166 righe). Insieme
sono un terzo del codice di tutti i moduli.

Che siano due, e che siano intercambiabili, \`e la prova che l'astrazione
introdotta nel \cref{cap:07-modalita-esecuzione} non \`e speculativa: il progetto
la esercita con due realizzazioni sostanzialmente diverse.

\section{\code{k8s-deployment-provider}}

Il provider Kubernetes traduce una \code{FunctionSpec} in tre oggetti: un
\code{Deployment}, un \code{Service} e --- se la strategia di scaling \`e
\code{HPA} --- un \code{HorizontalPodAutoscaler}. Usa il client
Fabric8.

\subsection{Convenzioni di denominazione ed etichette}

I nomi degli oggetti sono derivati dal nome della funzione con prefisso
\code{fn-}, e ogni oggetto porta due etichette: \code{app=nanofaas} e
\code{function=<nome>}. La seconda \`e anche il selettore del Deployment.

Il template dei pod porta tre annotazioni per lo scraping Prometheus:

\begin{lstlisting}[language=Java]
.addToAnnotations(NanofaasDeploymentConstants.ANNOTATION_SCRAPE, "true")
.addToAnnotations(NanofaasDeploymentConstants.ANNOTATION_PATH, "/metrics")
.addToAnnotations(NanofaasDeploymentConstants.ANNOTATION_PORT, "8080")
\end{lstlisting}

Le metriche lato funzione sono quindi raccolte senza che l'operatore configuri
nulla. \`E una piccola decisione con una conseguenza sperimentale grande: le
misure del \cref{cap:25-valutazione-concorrenza} che confrontano il tempo di
servizio misurato dal control plane con quello misurato dentro la funzione sono
possibili solo perch\'e la seconda sorgente esiste per default.

\subsection{Le variabili d'ambiente riservate}

\begin{lstlisting}[language=Java]
private static final Set<String> RESERVED_ENV = Set.of(
        "FUNCTION_NAME", "WARM", "TIMEOUT_MS", "EXECUTION_MODE", "WATCHDOG_CMD", "CALLBACK_URL"
);
\end{lstlisting}

Sono il contratto fra control plane e runtime della funzione
(\cref{cap:17-runtime-funzione}). Il builder le inietta e impedisce alla
specifica dell'utente di sovrascriverle: \code{CALLBACK\_URL} in particolare \`e
l'indirizzo a cui il runtime invia il completamento asincrono, e un valore
sbagliato produrrebbe esecuzioni che non terminano mai --- il caso che il terzo
orizzonte di eviction dello \code{ExecutionStore} deve poi ripulire
(\cref{cap:06-nucleo}).

\subsection{Il deadlock dello scale-to-zero con HPA}

Il commento pi\`u istruttivo del modulo:

\begin{lstlisting}[language=Java]
if (replicas == 0 && hpaOwnsScaling(spec.scalingConfig())) {
    // Kubernetes 1.36 only lets an HPA scale a target up from zero once it
    // has parked that target at zero itself (status condition ScaledToZero).
    // A Deployment created at zero therefore deadlocks: "scaling is disabled
    // since the replica count of the target is zero", forever. Start at one
    // and let the HPA scale it down, which is what stamps that condition.
    replicas = 1;
}
\end{lstlisting}

Una funzione con \code{minReplicas = 0} e strategia \code{HPA} verrebbe creata a
zero repliche, e l'HPA rifiuterebbe di risvegliarla perch\'e non \`e stato lui a
portarla a zero. Il rimedio \`e creare il Deployment a una replica e lasciare che
sia l'HPA a scalarla a zero, il che imprime la condizione di stato che poi gli
consente di risalire.

\`E il tipo di dettaglio che vale la pena riportare in un documento accademico
per una ragione metodologica: una piattaforma FaaS costruita sopra Kubernetes
eredita i vincoli di Kubernetes, e alcuni di questi non sono documentati come
vincoli ma si manifestano come comportamenti. Il costo dell'astrazione
``deployment gestito'' include l'obbligo di conoscerli.

\subsection{Validazione delle immagini}

\code{KubernetesImageValidator} verifica che un'immagine esista e sia scaricabile
\emph{prima} che la registrazione della funzione venga accettata. Il metodo \`e
diretto: crea un pod di prova, ne osserva lo stato del container, e cerca le
condizioni di attesa che segnalano un fallimento di pull
(\code{ErrImagePull}, \code{ImagePullBackOff}), con un timeout di 20 secondi e
polling a 500\,ms.

\`E un controllo costoso --- crea e distrugge un pod --- eseguito una sola volta
per registrazione. L'alternativa sarebbe scoprire l'immagine inesistente alla
prima invocazione, quando il chiamante sta gi\`a attendendo e la diagnosi \`e
molto meno diretta. La registrazione fallisce con un errore che nomina il
problema.

\section{\code{container-deployment-provider}}

Il provider locale crea container su un runtime compatibile con Docker. \`E
pensato per lo sviluppo e --- soprattutto --- per gli esperimenti su una singola
macchina, dove interporre Kubernetes fra il control plane e le funzioni
aggiungerebbe una variabile senza aggiungere realismo.

\subsection{Struttura}

\begin{center}
\small
\begin{adjustbox}{max width=\textwidth}
\begin{tabular}{@{}lp{8.6cm}@{}}
\toprule
\textbf{Componente} & \textbf{Ruolo} \\
\midrule
\code{ContainerRuntimeAdapter} & Astrazione sul runtime; due implementazioni,
via CLI e via client Docker Java. \\
\code{EphemeralPortAllocator} & Assegna una porta libera dell'host a ciascun
container. \\
\code{HttpEndpointProbe} & Attende che il container risponda sull'endpoint di
salute. \\
\code{RoundRobinFunctionProxy} & Espone un endpoint stabile che bilancia sulle
repliche. \\
\code{DockerImageValidator} & Verifica l'esistenza dell'immagine. \\
\bottomrule
\end{tabular}
\end{adjustbox}
\end{center}

\subsection{Il proxy: perch\'e serve}

Kubernetes fornisce un \code{Service}, ossia un nome stabile che bilancia sui
pod. Un runtime container locale non ha nulla di equivalente: ogni container ha
la propria porta effimera sull'host. Se il control plane conoscesse quelle porte
direttamente, ogni scale-up e ogni scale-down cambierebbero l'endpoint della
funzione.

\code{RoundRobinFunctionProxy} colma la differenza con un piccolo server HTTP ---
costruito sul \code{HttpServer} del JDK, senza dipendenze aggiuntive --- che
espone \code{/invoke} e \code{/health} su una porta stabile e distribuisce le
richieste sui backend correnti con round robin. L'insieme dei backend \`e un
riferimento atomico che il provider aggiorna a ogni cambio di scala.

\`E, in sostanza, una reimplementazione minimale del \code{Service} di
Kubernetes, ed \`e la componente che rende i due provider \emph{equivalenti dal
punto di vista del control plane}: entrambi restituiscono un URL stabile.

\subsection{Il \code{cpuset}: una decisione derivata da una misura}

Il commento su \code{ContainerLocalProperties} \`e uno dei pi\`u densi del
progetto:

\begin{quote}\itshape
\code{cpuset}: le CPU dell'host a cui ogni container di funzione gestita \`e
vincolato, per esempio \code{"0-3"}. Un limite di CPU per funzione limita ciascun
container separatamente, il che su una macchina con core liberi significa che le
funzioni non competono mai davvero --- misurato: due funzioni con quattro CPU
ciascuna su un host a undici core hanno condiviso un control plane e si sono a
malapena influenzate a vicenda. Vincolarle agli stessi core \`e ci\`o che rende la
capacit\`a della piattaforma una quantit\`a che devono dividersi, che \`e la
premessa su cui \`e costruito il budget di concorrenza.
\end{quote}

Questa \`e una funzionalit\`a nata da un requisito \emph{sperimentale}, non
operativo. Il modo \code{BUDGETED} del \cref{cap:12-concurrency-control} presume
che la concorrenza sia una risorsa scarsa condivisa; su un host con core in
eccesso quella premessa \`e falsa e il modo non pu\`o essere valutato. Il
\code{cpuset} fabbrica la scarsit\`a necessaria a rendere l'esperimento
significativo.

\`E anche il parametro che ha consentito di \emph{escludere} la contesa di CPU
come spiegazione del cross-talk fra funzioni co-residenti: vincolando entrambe
agli stessi quattro core, i risultati sono rimasti identici a tre cifre
significative (\cref{cap:25-valutazione-concorrenza}).

\subsection{Readiness}

\code{HttpEndpointProbe} interroga l'endpoint di salute del container con un
timeout di 20 secondi e un intervallo di 250\,ms. Un container che non diventa
pronto entro il timeout fa fallire il provisioning.

I valori sono deliberatamente pi\`u aggressivi delle sonde di Kubernetes:
l'obiettivo qui non \`e la robustezza in produzione ma il tempo di ciclo di un
esperimento, in cui il provisioning avviene decine di volte.

\section{Chi possiede la validazione delle immagini}

Un punto architetturale gi\`a anticipato nel \cref{cap:03-requisiti-principi} e
che qui trova la sua giustificazione completa: la validazione delle immagini
\emph{non} \`e un modulo autonomo. Ciascun provider possiede il proprio
validatore e lo attiva quando \`e selezionato come backend.

Il motivo \`e che l'operazione non \`e la stessa. Verificare un'immagine tramite
Kubernetes significa creare un pod e osservare le condizioni di attesa del
container; verificarla tramite un demone Docker locale significa interrogarne
l'indice delle immagini o tentare un pull. Un modulo generico avrebbe dovuto
reintrodurre al proprio interno la distinzione fra backend che l'astrazione dei
provider esiste per eliminare, e avrebbe dovuto essere selezionato
coerentemente con essi --- un vincolo di coerenza in pi\`u nella selezione
modulare, per nessun guadagno.

Resta il fatto che l'interfaccia \code{ImageValidator} vive nel package
\code{registry} del control plane anzich\'e in \code{platform/common}, dove
stanno gli altri contratti condivisi (\cref{cap:05-contratti-dati}). Non \`e una
violazione architetturale --- i moduli possono dipendere dal nucleo, e la regola
ArchUnit vieta soltanto le dipendenze \emph{fra} moduli
(\cref{cap:26-qualita-verifica}) --- ma \`e una collocazione incoerente con il
criterio dichiarato altrove, e il posto naturale dell'interfaccia \`e il modulo
dei contratti.

\section{Confronto}

\begin{table}[htbp]
\centering
\caption{I due provider di deployment a confronto.}
\label{tab:provider}
\small
\begin{tabularx}{\textwidth}{@{}lXX@{}}
\toprule
 & \textbf{\code{k8s}} & \textbf{\code{container-local}} \\
\midrule
Backend & API Kubernetes (Fabric8) & runtime Docker-compatibile \\
Endpoint stabile & \code{Service} & proxy round robin interno \\
Scaling & \code{replicas} del Deployment o HPA & creazione/distruzione
container \\
Validazione immagine & pod di prova, condizioni di attesa & indice del demone \\
Isolamento CPU & \code{resources.limits} & \code{resources.limits} +
\code{cpuset} \\
Dipendenze pesanti & Fabric8, Vert.x & client Docker Java \\
Uso previsto & carichi realistici su cluster & esperimenti su singola macchina \\
\bottomrule
\end{tabularx}
\end{table}

L'ultima riga \`e la ragione per cui entrambi esistono. Il provider Kubernetes
\`e quello che rende \nf{} una piattaforma FaaS reale; quello locale \`e quello
che rende gli esperimenti eseguibili in un ciclo abbastanza breve da poterli
ripetere.
